% Generated draft: requires mathematical and source review.
\documentclass[11pt]{article}
\usepackage[margin=25mm]{geometry}
\usepackage{fontspec}
\setmainfont{Latin Modern Roman}
\usepackage{amsmath,amssymb,mathtools,unicode-math}
\setmathfont{Latin Modern Math}
\usepackage{enumitem,needspace,xurl,hyperref}
\setlength{\parindent}{0pt}
\setlength{\parskip}{5pt}
\setlength{\emergencystretch}{3em}
\newcommand{\N}{\mathbb{N}}
\newcommand{\Z}{\mathbb{Z}}
\newcommand{\Q}{\mathbb{Q}}
\newcommand{\R}{\mathbb{R}}
\newcommand{\C}{\mathbb{C}}
\newcommand{\F}{\mathbb{F}}
\newcommand{\A}{\mathbb{A}}
\newcommand{\PP}{\mathbb{P}}
\newcommand{\E}{\mathbb{E}}
\newcommand{\eps}{\varepsilon}
\newcommand{\dd}{\,\mathrm{d}}
\begin{document}
\section*{UnsolvedMath ID: 30000644}
\textbf{Source problem code:} catalogue entry
\subsection{Definitions and mathematical statement}
Let \(R\) be a commutative unital ring containing \(\Q\), let \(n,m\in\Z_{>0}\), and let \(t,y_1,\ldots,y_n\) be commuting indeterminates. Put \(A=R[t]\), \(\overline A=R[t]/(t^m)\), and \(Y=(y_1,\ldots,y_n)\). For a commutative ring \(B\), define
\[
 \operatorname{SAut}_{B}(B[Y])=
 \left\{F\in\operatorname{Aut}_{B}(B[Y]):
 \det\left(\frac{\partial F_i}{\partial y_j}\right)_{i,j=1}^n=1\right\},
\]
where \(F_i=F(y_i)\) and the determinant equality is in \(B[Y]\). Thus \(F\) must already be an algebra automorphism, not just a polynomial endomorphism with unit Jacobian. Its inverse is also polynomial. Differentiation is only in the \(Y\)-variables; all base elements, including \(t\), are fixed.

Coefficient reduction modulo \(t^m\) induces
\[
 \rho_{R,n,m}:\operatorname{SAut}_{R[t]}(R[t][Y])
 \longrightarrow
 \operatorname{SAut}_{\overline A}(\overline A[Y]).
\]
Reducing an automorphism and its inverse preserves the inverse identities, and the Jacobian determinant reduces to one, so this is the stated map.

The question is whether \(\rho_{R,n,m}\) is surjective for every stated \(R,n,m\). Explicitly, for each \(\overline F\) in the target, must there exist \(F\) in the source such that
\[
 F_i\bmod t^m=\overline F_i\quad\text{for }1\leq i\leq n?
\]
The lift must have determinant exactly one in \(R[t][Y]\), rather than just congruent to one, and must be a genuine polynomial automorphism over \(R[t]\). A formal \(t\)-adic lift, a polynomial endomorphism without an inverse, or a lift after adding variables would not answer the question. There is no bound on the degrees of a lift in the supplied wording.

No domain, field, reducedness, Noetherianity or finite-generation hypothesis on \(R\) is inserted; the source explicitly permits a general commutative ring containing \(\Q\). The number of variables remains \(n\). The literature triage checked 2026-08-21 reports no authoritative complete resolution located. The open label and notation-only verification are dated source metadata, not an independently checked current theorem.
\subsection{Short English statement}
Must every special polynomial automorphism over the truncated coefficient ring lift to a special polynomial automorphism over \(R[t]\), for every commutative \(\Q\)-containing ring? The lift must be polynomially invertible, preserve the same variables and have Jacobian exactly one. Formal or stabilized lifting is a different assertion.
\subsection{Sources}
\begin{itemize}
\item Local UnsolvedMath registry, permanent ID 30000644, source code catalogue entry, ``Surjectivity of Reduction Maps for Special Polynomial Automorphisms,'' including original and cleaned map, special-automorphism definition and verification metadata.
\item \emph{Affine Algebraic Geometry (2007)}, Oberwolfach Reports, pp.\ 5--82, DOI 10.4171/owr/2007/01: \url{https://doi.org/10.4171/owr/2007/01}. This is also the sole supplied literature reference.
\item Literature triage checked 2026-08-21; local record updated 2026-09-16. No external source has been visited or independently checked here.
\end{itemize}
\end{document}
